% Source: https://artofproblemsolving.com/community/c3h2652398_2017_amc8_25
% Stable identity: aops-topic-2652398-q01
A one-sixth arc of radius $2$ subtends $60^\circ$ at its centre, so its chord has length $2$: $TR=SR=2$. Also $TS=2$ (triangle $UTS$ is isosceles with a $60^\circ$ apex, hence equilateral). So $R$ is the reflection of $U$ in $TS$, and $UTRS$ is a rhombus made of two equilateral triangles of side $2$, with area $2\cdot\sqrt3=2\sqrt3$.

Each arc bends into the rhombus, cutting off a circular segment between the arc and its chord. A segment of a $60^\circ$ arc of radius $2$ has area
\[\dfrac16\pi\cdot2^2-\dfrac{\sqrt3}{4}\cdot2^2=\dfrac{2\pi}{3}-\sqrt3.\]
The region is the rhombus minus two such segments:
\[2\sqrt3-2\left(\dfrac{2\pi}3-\sqrt3\right)=\boxed{4\sqrt3-\dfrac{4\pi}{3}}.\]

\begin{asy}
usepackage("amsmath");
size(3.5cm);
defaultpen(fontsize(10pt));
pair U=(0,0),T=2*dir(120),Sx=2*dir(60);pair R=T+Sx;
pair X1=T+R-Sx,X2=Sx+R-T;
fill(U--T--arc(X1,2,-60,0)--arc(X2,2,180,240)--cycle,lightgray);
draw(U--T,linewidth(0.8));draw(U--Sx,linewidth(0.8));
draw(arc(X1,2,-60,0),linewidth(0.8));draw(arc(X2,2,180,240),linewidth(0.8));
draw(T--R--Sx--cycle,dashed);draw(X1--T,dotted);draw(X1--R,dotted);draw(X2--R,dotted);draw(X2--Sx,dotted);
dot(X1);dot(X2);
label("$U$",U,S);label("$T$",T,W);label("$S$",Sx,E);label("$R$",R,N);label("$X$",X1,NW);label("$Y$",X2,NE);
\end{asy}

(The arcs have centres $X$ and $Y$, which form equilateral triangles $TRX$ and $SRY$ outside the rhombus.)
